\documentclass[12pt]{article}
\usepackage[a4paper,margin=2cm]{geometry}
\usepackage{amsmath,amssymb}
\usepackage{xepersian}
\settextfont[
  Path=../../Homework/1401-2/HW2 - official solution/,
  Extension=.ttf,
  BoldFont=*Bd,
  ItalicFont=*It,
  BoldItalicFont=*BdIt
]{XB Niloofar}
\setdigitfont[
  Path=../../Homework/1401-2/HW2 - official solution/,
  Extension=.ttf,
  BoldFont=*Bd,
  ItalicFont=*It,
  BoldItalicFont=*BdIt
]{XB Niloofar}
\setlength{\parindent}{0pt}
\setlength{\parskip}{0.55em}

\begin{document}
\begin{center}
  {\Large\textbf{اصلاحیهٔ پاسخ رسمی پایان‌ترم ـ نیم‌سال ۱۴۰۲--۱}}
\end{center}

این برگه پاسخ دو سؤال تشریحی را که محیط پاسخ آن‌ها در فایل \lr{TeX} منبع خالی مانده است تکمیل می‌کند و باید همراه فایل اصلی خوانده شود.

\section*{بخش تشریحی، سؤال ۳}
با فرض آن‌که تکرار در تعداد متناهی گام در
\(x^*\)
متوقف نشود، خطا را با
\(e_k=x^k-x^*\)
نشان می‌دهیم. از تیلور و فرض
\[
 g^{(j)}(x^*)=0,\qquad j=1,\ldots,p-1
\]
برای عددی \(\xi_k\) میان \(x^k\) و \(x^*\) داریم
\[
  e_{k+1}
  =g(x^k)-g(x^*)
  =\frac{g^{(p)}(\xi_k)}{p!}e_k^p.
\]
چون \(x^k\to x^*\)، داریم \(\xi_k\to x^*\). پیوستگی \(g^{(p)}\) در همسایگی \(x^*\) نشان می‌دهد این مشتق روی یک بازه‌ی بسته‌ی کوچک‌تر پیرامون \(x^*\) کراندار است؛ پس برای ثابت \(\bar c\)
\[
  |e_{k+1}|\leq \bar c\,|e_k|^p
\]
و مرتبه دست‌کم \(p\) است. افزون بر آن،
\[
  \lim_{k\to\infty}
  \frac{|e_{k+1}|}{|e_k|^p}
  =\frac{|g^{(p)}(x^*)|}{p!}.
\]
اگر \(g^{(p)}(x^*)\neq0\)، این حد مثبت و متناهی است و مرتبه دقیقاً \(p\) خواهد بود. اگر این شرط حذف شود، فقط نتیجهٔ «دست‌کم \(p\)» تضمین می‌شود.

\section*{بخش تشریحی، سؤال ۲}
در صورت منتشرشده، ضریب جمله‌ی دوم نیز \(c_1\) نوشته شده است؛ این یک خطای تایپی است و باید \(c_2\) باشد. بنابراین از
\[
 I=F_1(h)+c_1h^p+\mathcal O(h^{p+1}),
 \qquad
 I=F_2(h)+c_2h^p+\mathcal O(h^{p+1}),
\]
با \(c_1\neq c_2\)، جملهٔ پیشرو خطا در ترکیب
\[
  \boxed{\;
  F_{\mathrm{new}}(h)
  =\frac{c_2F_1(h)-c_1F_2(h)}{c_2-c_1}
  \;}
\]
حذف می‌شود.

در بخش ب، منظور از قاعدهٔ مستطیلی باید \textbf{قاعدهٔ نقطهٔ میانی} باشد. با
\[
 T=\frac h2\bigl(f(a)+f(b)\bigr),
 \qquad
 M=h\,f\!\left(\frac{a+b}{2}\right),
 \qquad h=b-a,
\]
بسط حول نقطهٔ میانی نشان می‌دهد جمله‌های پیشرو خطای \(T\) و \(M\) نسبت
\(-2\)
دارند؛ بنابراین
\[
  S=\frac{T+2M}{3}
   =\frac h6
     \left[
       f(a)+4f\!\left(\frac{a+b}{2}\right)+f(b)
     \right],
\]
که همان قاعدهٔ سیمپسون است. برای توجیه افزایش مرتبه و کران خطای معمول سیمپسون کافی است
\(f\in C^4[a,b]\)
باشد؛ در این صورت خطای قاعدهٔ تک‌بازه‌ای از مرتبهٔ
\(\mathcal O(h^5)\)
است.
\end{document}
